\begin{homeworkProblem}
  Seja $H$ um subconjunto não vazio de um grupo $G$. Mostre que $H \le G$ se, e somente se, para todo $a,b \in H$ tem-se $ab^{-1} \in H$.
  \begin{solution}
    (\textbf{Suficiente}) Si $H\leq G$, então para tudo $a,b\in H$ tem-se $ab^{-1}\in H$.

    Como $H$ é um subgrupo, em particular $H$ é grupo, então si $a,b\in H$ ten-se por a propiedade do existencia do elemento inverso que $b^{-1}\in H$, logo por a propiedade do fechamento $ab^{-1}\in H$.\\

    (\textbf{Necessária}) Si para tudo $a,b\in H$ tem-se $ab^{-1}\in H$, então $H\leq G$.
    \begin{itemize}
      \item Existencia do elemento neutro.

        Como $a\in H$, então por hipótese $aa^{-1}=1\in H$.
      \item Existencia do elemento inverso.

        Como $1,a\in H$, então $1a^{-1}=a^{-1}\in H$.
      \item Fechamento.

        Como $a,b\in H$, por existencia do elemento inverso $b^{-1}\in H$, logo usando la hipótese con $a$ e $b^{-1}$ se sabe que $ab\in H$.
      \item Associatividade.

        Isso ocorre porque $G$ é um grupo.
    \end{itemize}
    Portanto, podemos concluir que a afirmação é verdadeira.
  \end{solution}
\end{homeworkProblem}
